\documentclass{article}
\usepackage[utf8]{inputenc}
\usepackage{amsmath, amsthm, amsfonts,amssymb}
\usepackage[spanish]{babel}
\usepackage{multicol}
\usepackage{listings}
\lstset{basicstyle=\footnotesize\ttfamily,breaklines=true}
\usepackage{alltt}
\usepackage{graphicx}
\usepackage{subfigure}
\usepackage{subfig}
\usepackage{float}
\usepackage{url}
\usepackage{enumerate}
\usepackage{framed}
\usepackage{color}
\usepackage{cancel}
\usepackage{wrapfig}\definecolor{shadecolor}{RGB}{250,250,250}
\usepackage{framed}
\usepackage{epstopdf}
\setlength\parindent{0pt}
\usepackage{listings}
\usepackage{color} %red, green, blue, yellow, cyan, magenta, black, white
% Operadores matemáticos y simbolos
\DeclareMathOperator{\dive}{div}
\DeclareMathOperator{\trace}{trace}
\DeclareMathOperator{\tr}{tr}
\DeclareMathOperator{\symm}{symm}
\DeclareMathOperator{\sk}{skew}
\DeclareMathOperator{\grad}{grad}
\DeclareMathOperator{\Grad}{Grad}
\DeclareMathOperator{\curl}{curl}
\DeclareMathOperator{\Curl}{Curl}
\def\R{\mbox{\(\mathbb{R}\)}}
\def\E{\mbox{\(\mathbb{E}\)}}
\def\P{\mbox{\(\mathbb{P}\)}}
\def\I{\mbox{\(\mathbb{I}\)}}
\def\L{\mbox{\(\mathbb{L}\)}}
\def\dx{\mbox{\(\,\mathrm{d}x\)}}
\usepackage{geometry}
\geometry{left=2.5cm, right=2.5cm, top=2cm, bottom=3cm}
\title{Tarea 4\\}
\author{Luis Felipe Silva De Vidts}
\begin{document}
\begin{figure}
\begin{minipage}{2.5cm}
\includegraphics[width=0.8\textwidth]{./figures/LogoUC-BN}
\end{minipage}
\begin{minipage}{14.5cm}
\vspace{4mm}
{\sc PONTIFICIA UNIVERSIDAD CAT\'OLICA DE CHILE}\\
Facultad de Matemáticas\\
Departamento de Estadísica\\
{\bf EYP1113 Probabilidades y Estadística}\\
\vspace{0mm}
\hrulefill
\end{minipage}
\end{figure}
\phantom{""}
\vspace{-5mm}
\normalsize
\begin{center}
\Huge Tarea 1\\
\normalsize Luis Felipe Silva De Vidts
\end{center}
\section*{Parte I}
% Aqui hay q poner la tabla con los estimados  al hacer el QQ-plot y la recta de ajuste
% Junto con los graficos de probabilidad asociados

\section*{Parte II}
\begin{enumerate}[a)]
\item 40 % Describa a partir de medidas empı́ricas y gráficos adecuados las variables: Latitud, Longitud, Magnitud, Profundidad, Radiacion, Presion, Humedad, Temperatura y VelocidadViento, que se encuentran en la Datos02.txt.
\item 20% Represente el comportamiento del número esperado de sismos por hora mediante una función simple con respecto al tiempo u otra variables, desde la ocurrencia del terremoto 6.9 o Richter e ilustre gráficamente. La función lm() y nls() son útiles para estimar coeficientes de funciones lineales y no lineales. El conteo de sismos se encuentra en Datos03.txt.
\end{enumerate}

\end{document}